% Section 5: Discussion, design implications, limitations, ethics and positionality, future work.
% Every number is a \jm... macro from make_macros.py. No hand-typed numbers.
% Cross-reference labels expected from other sections: sec:default, sec:guard, sec:waterfall, sec:validity, sec:frontier, sec:req, sec:design.
\section{Discussion}\label{sec:discussion}

\subsection{What the evidence establishes}\label{sec:disc-summary}

Three results carry the paper, and one measurement result qualifies how far they reach.
First, the affirmative default is a property of the model family, not of one prompt: with bare keywords every one of the \jmbreadthModels{} models (\jmbreadthFamilies{} families) chooses the affirmative reading in at least \jmbreadthMinAff{}\% of inputs, and without the prompt legend the pooled share is still \jmlegendNoLegendAff{}\% [\jmlegendNoLegendAffLo{}, \jmlegendNoLegendAffHi{}].
We do not present the direction of this bias as new; it builds on the positivity bias that Chen et al.\ report for keyword-to-sentence generation about negative commonsense facts~\cite{chen2023say}, adds the base rate in human-written keyword messages, and prices the default per dimension.
At the human base rate of negation in keyword messages (\jmtTwoNegBase{}\%), always committing to the default leaves \jmfOneRevPolBare{} silent polarity reversals per 1,000 messages.
Second, a one-token cue removes most of that cost: with the cue on every message, polarity reversals fall to \jmfOneRevPolChip{} (chip) or \jmfOneRevPolWord{} (plain word) per 1,000, politeness reversals from \jmfOneRevPoliteBare{} to \jmfOneRevPoliteChip{} and \jmfOneRevPoliteWord{}, and chip plus guard raises speech-act restoration from \jmfOneActRestoreChip{}\% to \jmfOneActRestoreChipG{}\%.
An irrelevant control cue does not do this (Section~\ref{sec:default}).
Third, the guard that repairs a failed draft by deletion raises speech-act restoration on the same items from \jmfOneActRestoreChip{}\% to \jmfOneActRestoreChipG{}\% (its transfer to fresh families is restated from our earlier paper in Section~\ref{sec:guard}), and the blind repair audit was inconclusive on usability: the difference was \jmpOneSevenDiffUsable{} points with interval \jmpOneSevenDiffUsableLo{} to \jmpOneSevenDiffUsableHi{}, which neither shows nor excludes a cost (Section~\ref{sec:guard}).
Against each repair's own failed original, the audit found large gains in preserved meaning (\jmpOneSevenOrigPres{} points) and fewer invented stances (\jmpOneSevenOrigStance{} points).

The qualifying result is that these gains concern a floor, not the target.
Writers accepted \jmpOneEightAcceptYesK{} of \jmpOneEightAcceptN{} drafts as written, and among drafts whose critical commitment the judge had marked correct, \jmwmRightROne{}\% were still not the version the writer intended (\jmwmRightROneK{} of \jmwmRightROneN{}).
The reasons were compression loss and misresolved events or roles, not added content: in \jmwmDroppedUnrecShare{}\% of the rejected cases the writer's sentence contained words absent from the keywords, so no model could have recovered them, and \jmwmRightSubst{}\% of the judge-correct rejections were event or role substitutions (who does what to whom).
Added content accounted for \jmwmAdded{}\% of primary reasons.
A cue nevertheless lifts the match to the writer's intended version from \jmpOneEightMatchBare{}\% to \jmpOneEightMatchChip{}\% (\jmpOneEightCueGain{} points, interval \jmpOneEightCueGainLo{} to \jmpOneEightCueGainHi{}, clustered over \jmpOneEightWriters{} writers), so the cue is the most effective single input measured here while remaining far from sufficient.
We therefore read critical-commitment preservation as a hard filter that a draft must pass before it is offered, and the coverage of the writer's whole intended message, per tap and per draft read, as the objective that remains open (Section~\ref{sec:future}).

\subsection{Design implications for AAC}\label{sec:implications}

The implications below follow from the measured results and from the requirement changes that two paid design consultants who use AAC proposed after reviewing the requirements and stored example outputs (Section~\ref{sec:req-revisions}).
We report the consultants' input as changes to the requirement list; their experiences are not data and the list remains our proposal.
Table~\ref{tab:implications} states each implication with its evidence and its status.
An implication marked \emph{measured} rests on a result in this paper; one marked \emph{consultant} was proposed by the consultants and is adopted as a requirement; one marked \emph{untested} is a design hypothesis that the studies in Section~\ref{sec:future} are built to test.

\begin{table}[t]
\caption{Design implications, the evidence behind each, and its status. Statuses: measured (result in this paper), consultant (requirement proposed by the two design consultants), untested (hypothesis for future work).}
\label{tab:implications}
\small
\begin{tabular}{@{}p{0.04\linewidth}p{0.36\linewidth}p{0.42\linewidth}p{0.12\linewidth}@{}}
\toprule
 & Implication & Evidence & Status \\
\midrule
I1 & Offer commitment cues (negation, politeness, sentence type) as first-class controls, available as a chip or as an equivalent plain word. & Polarity reversals per 1,000: \jmfOneRevPolBare{} bare, \jmfOneRevPolChip{} chip, \jmfOneRevPolWord{} word; politeness \jmfOneRevPoliteBare{}, \jmfOneRevPoliteChip{}, \jmfOneRevPoliteWord{}. Chip and word are not distinguishable on polarity and are equivalent within the pre-specified margin on politeness. & measured \\
I2 & Do not fix one route as the default. Both routes set the same visible meaning and allow correction in place. & One consultant found the chip easier to imagine; the other answered that it depends on the access method and vocabulary layout. Neither claim is a measure of selection effort. & consultant; untested \\
I3 & A generated sentence is a draft. It is spoken only after the person approves it, can be inspected privately, edited or cancelled, and the original selection can always be spoken directly. & Proposed independently by both consultants as a sixth requirement (R6). Writers accepted \jmpOneEightAcceptYesK{} of \jmpOneEightAcceptN{} drafts as written, so unapproved speech would voice a draft the writer would change in nearly every case. & consultant; measured \\
I4 & Keep exact-copy as the default for fully typed sentences; rewriting needs an explicit request, a visible proposed change and approval. & Both consultants. Our studies cover keyword input only, so this is a scope boundary, not a result. & consultant \\
I5 & Where a draft must choose between readings, show a small number of meaningfully different drafts with the difference highlighted, and reserve a focused question for cases that block any useful draft. & Cue-supplied match to the intended version is \jmpOneEightMatchChip{}\%; the remainder is event, role and detail the cue does not carry. Asking at a fixed budget does not beat random choice of where to ask (I7). Both consultants asked for differing drafts or a targeted choice. & consultant; untested \\
I6 & Add event and role structure (who does what to whom) to the commitment model, alongside the seven dimensions. & \jmwmRightSubst{}\% of judge-correct rejections were event or role substitutions; both consultants independently asked that relations be preserved. & measured; consultant \\
I7 & Do not allocate effort to ask-when-uncertain policies. Supply the cue instead. & Ask policies do not beat gated-random asking at a 10\% budget: prior-ranked \jmdpAskPriorPol{} [\jmdpAskPriorPolLo{}, \jmdpAskPriorPolHi{}], gated-random \jmdpAskRandPol{} [\jmdpAskRandPolLo{}, \jmdpAskRandPolHi{}], a CLAM-style gate \jmclamPol{} [\jmclamPolLo{}, \jmclamPolHi{}] polarity reversals per 1,000, against \jmdpAlwaysCommitPol{} for always-commit. Minimum detectable difference \jmdpMdePol{} per 1,000. & measured (null of limited power) \\
I8 & Spend cue taps by clinical severity as well as frequency. & Two SLPs rated negation errors \jmslpNegSeverity{} of 5 and would block them (\jmslpNegBlock{}\% of ratings). Severity weights: polarity \jmwPol{}, speech act \jmwAct{}, politeness \jmwPolite{}. Speech act is \jmwShareAct{}\% of weighted failure under always-commit because it is frequent (detector-proxy labels, provisional). & measured (n=\jmdesignSlp{} experts) \\
\bottomrule
\end{tabular}
\end{table}

\paragraph{A concrete interface recommendation, and what is unmeasured.}
Our evidence supports a specific design: show one draft at a time, with the selected keywords and any cue visible beside it; make the first action ``speak the original selection'' and the second ``approve this draft'', never a default; offer the negation and politeness cues as a standing control that costs one selection before generation; and offer a second, differing draft only when the cue has been supplied and the writer rejects the first. We rank this design on meaning restored and acceptance, not on tap or reading cost: the paper never measured taps, reading time, scanning order or gaze effort, so ``lowest cost'' is the objective of the study that follows and not a result. A cue costs the writer a selection in a direct-touch layout, and may cost far more under scanning.

Two points in the table deserve emphasis.
The first is the approval gate (I3).
A keyword system that speaks its best sentence converts every residual error into a statement the AAC user did not choose, and the writer-grounded numbers show the residual is large even when the critical commitment survives.
The consultants framed the gate as control over delivery, not as a safety feature: drafting and editing never speak or send, the person decides when a draft is spoken, and the person may keep an unusual, blunt, short or incomplete sentence.
This also constrains the guard.
A deletion-based repair is acceptable as a proposal that the person sees and approves; it is not acceptable as a silent edit of what is spoken.
The SLPs reflect the same caution from the clinical side: they judged \jmslpAccRepair{}\% of repairs acceptable to speak as they stand (\jmslpAccRepairN{} rater-judgments), against \jmslpAccOrig{}\% of the originals for the same inputs.

The second is the choice between asking and cueing (I5, I7).
Ask policies are attractive because a question costs one turn.
Our comparison shows no resolvable advantage of any tested policy over gated-random asking at the same budget, and in our machine-only checks the model's own uncertainty did not reliably mark the inputs at risk.
Giving the writer the cue as a standing control is the efficient option on the dimensions we measured.
Reading cost is the unmeasured side of this trade: a set of differing drafts replaces taps with reading, and whether writers who use scanning, switches or gaze can catch a wrong draft in a set is untested.
The approval gate presupposes that they can.

\paragraph{Implications for evaluation practice.}
The measuring instrument is itself a result.
Automatic judges agreed with a human arbiter on \jmpOneNineAgreeJOne{}\% (judge 1) and \jmpOneNineAgreeJTwo{}\% (judge 2) of \jmpOneNineCodings{} codings, against \jmpOneNineAgreeArb{}\% (\ensuremath{\kappa}=\jmpOneNineKappaArb{}) for human coders, and their agreement with the writers' own judgment was low (\ensuremath{\kappa}=\jmpOneEightKappa{}) as was their agreement with the SLPs (\ensuremath{\kappa}=\jmslpJudgeKappaOne{} and \jmslpJudgeKappaTwo{}).
Whether ``Can you\ldots?'' is a request or a question is a convention that a judge can code differently from people: \jmpOneNineCanYouRequest{}\% of human codings (\jmpOneNineCanYouK{} of \jmpOneNineCanYouN{}) read it as a request.
We draw three practices from this.
Systems that claim meaning preservation should report agreement with a human arbiter or with the writer, not only with a model judge, in line with the argument that AI-powered AAC needs technical, user-mediated and functional evaluation together~\cite{frisch2026complicated}; known judge biases~\cite{zheng2023judging} make this a requirement, not a preference.
Failure rates should be reported unweighted and severity-weighted, because the weights state what the rate means clinically: the weighted cue failure falls from \jmslpWCueBare{}\% to \jmslpWCueChip{}\% (chip) and \jmslpWCueWord{}\% (word), against unweighted \jmslpUCueBare{}\% to \jmslpUCueChip{}\% and \jmslpUCueWord{}\%, and the weighted guard failure falls from \jmslpWGuardBase{}\% to \jmslpWGuardG{}\% (unweighted \jmslpUGuardBase{}\% to \jmslpUGuardG{}\%).
And a headline figure should name its denominator and its assumption: the ask-policy arithmetic assumes that the asked person supplies the cue.

\subsection{Limitations}\label{sec:limitations}

We state each limitation once, with the claim it bounds.

\paragraph{Who contributed the human data.}
Writers, annotators and raters were adult non-author university classmates, consented and paid (payment and food).
They are not AAC users, and the writer-intent results (acceptance, intended-version match) are agreement between writers and drafts built from their own keywords, not validation by AAC users.
We make no claim of benefit to AAC users and no claim about how AAC users would write keywords, which cues they would use or how they would review drafts.
The \jmdesignSlp{} speech-language pathologists are an expert review of a fixed set of outputs, not a population estimate, and their agreement statistics describe this pair.
The \jmdesignConsultants{} consultants are not a sample of AAC users; their feedback shaped the requirements, concerned a text-based design only, and does not establish what AAC users in general want.

\paragraph{Scale and status of the replication.}
The human-written pair replication has \jmhwChipGainWriters{} to \jmhwAffShareWriters{} writers per contrast and is exploratory; we claim direction and consistency across writers and models (the chip gain over a non-critical cue is \jmhwChipGain{} points [\jmhwChipGainLo{}, \jmhwChipGainHi{}] over \jmhwChipGainFamilies{} families), not interval width.
The study is pre-specified and git-timestamped; it is not pre-registered.

\paragraph{The usability audit.}
The interval includes zero but is wide: the lower bound is \jmpOneSevenDiffUsableLo{} points, and the repaired texts are less often usable and preserving together than controls (\jmpOneSevenDiffUsPres{} points).
Control items are easier, so the gap is confounded with item difficulty and we do not read it as a cost of repair.
We therefore describe the audit as inconclusive, with the interval shown, and not as ``no cost'' or ``as good as controls''.
Repairs of the smallest local model are usable in \jmpOneSevenGemmaUsable{}\% of cases (\jmpOneSevenGemmaUsableK{} of \jmpOneSevenGemmaUsableN{}), which bounds where the guard can be recommended.
Each item was rated by \jmdesignRatersPerItem{} independent raters, all classmates.

\paragraph{Measurement of the whole message.}
The whole-message analysis is post hoc.
The reason classifier agrees with a hand coder at $\kappa=\jmwmHandKappa{}$ on a \jmwmHandN{}-item check; because the hand coder is an LLM and the classifier rules were written after seeing those items, the figure is optimistic and not held out.
Compression loss is also partly a property of keyword input: when a writer's sentence carries content that the keywords never contained, no single draft can be accepted as written.
The acceptance rate (\jmpOneEightAcceptYesK{} of \jmpOneEightAcceptN{}) is therefore best read as a ceiling on what the single-draft paradigm can reach, not as a ranking of the four models rated.

\paragraph{Cost model and frontier.}
The base rate of negation (\jmtTwoNegBase{}\% of keyword messages) comes from a posterior over a labelled corpus; speech-act and politeness rates rest on detector-proxy labels and are marked provisional.
The ask-policy comparison is a null of limited power (minimum detectable difference \jmdpMdePol{} per 1,000 for polarity), not an equivalence result.
The reversals-per-1,000 arithmetic assumes the asked person supplies the cue; the same-items measurement on \jmfOneUnits{} units replaces the earlier transported estimate of that quantity.

\paragraph{Scope of the system studied.}
All inputs are English text keyword strings; we did not study speech output, access methods, vocabulary layout, prediction interfaces, multi-turn context or sustained use.
Model outputs come from specific model versions and a hosted API that does not honour a random seed, so exact outputs are not reproducible even though the stored outputs, judgments and analyses are.
% SLOT: add one sentence each for the pending human checks (treated outputs, source labels) and the intent-coverage pilot once their results exist.
Until those human checks return, the treated-output and source-label results above rest on machine labels.

\subsection{Ethics and positionality}\label{sec:ethics-positionality}

Annotators were adult non-author university classmates who gave consent and received payment and food.
They labelled and audited model outputs; some also chose keywords and wrote sentence pairs that we release under CC~BY~4.0 with their consent.
We collected no personal data, no demographic, health or disability information, and no recordings; participants are identified only by an ID, could stop at any time and were paid for the time spent.
Two speech-language pathologists were paid expert annotators of study materials, working anonymously under a written professional-services agreement; the data concern the outputs, not the raters.
Two paid design consultants who use AAC were engaged who reviewed requirements and examples under a written consulting agreement. This was design consultation, not research participation, and we (the authors) decided that it was; we did not collect, analyse or report their experiences, and we report only which requirements changed and why.
As independent researchers we did not seek an institutional review determination, and we make no claim of approval or exemption.
A study that recruits AAC users to use or evaluate a system requires institutional review before recruitment, and none is reported here.

The main ethical risk this work addresses is also its principal limit.
A system that expands a person's words into a fluent sentence takes part in authorship: when it is wrong it does not signal an error, it speaks for the person, and the first-person voice makes that harm easy to miss.
Prior work on AI in AAC records that language models can both help and hinder communication~\cite{valencia2023less}, and our results give a quantified instance: a missing cue silently produces the opposite polarity or the wrong sentence type at a measurable rate, and a draft that keeps every critical commitment is still rarely the sentence the writer meant.
For this reason we argue for approval-gated drafts, exact-copy defaults and visible differences, and we present the guard as a proposal generator, not an editor of what is spoken.
The cue and guard results are measurements of model behaviour on classmates' keyword messages; they are not evidence that deploying either would benefit a person who relies on AAC, and the paper should not be cited for that.
The released pairs and stored outputs contain no personal data, but they could be used to tune a system to pass our checks without improving communication; we state the checks as a floor for this reason.

Positionality.
This paper takes up the design of tools for a community that the data collected here do not represent.
Our consultants and SLPs shaped requirements and severity weights but did not review the analysis or approve the claims, and the requirements remain our proposals; we use the terms expert annotator and design consultant and avoid ``co-design'' and ``participatory'' because neither describes what happened.
% AUTHOR-INPUT: add a one- or two-sentence statement of the authors' own relationship to AAC, disability and language technology, if the authors choose to disclose one; do not invent it.
% AUTHOR-INPUT: insert the ACM generative-AI disclosure text if the venue requires one; the lab supplies it.

\subsection{Future work}\label{sec:future}

\paragraph{A study with AAC users.}
The central open question is whether the results transfer to the people for whom the tools are intended.
A study with AAC users, under institutional review and with an access-method-aware protocol, would test what this paper cannot: which cues users choose to supply and in which route; how often they can detect a wrong draft in a single draft and in a set; the reading and selection effort for scanning, switch, eye-gaze and direct-touch users; whether exact-copy and approval gates are used as the consultants expect; and whether users want the system to learn preferences from a single approval, which the consultants cautioned against.
The consultants also noted that the chip and plain-word routes need a working prototype with real navigation and scanning order to compare, because our evidence on the two routes concerns meaning restored, not effort.

\paragraph{Intent coverage of approval-gated draft sets.}
The measure we propose to replace single-draft correctness is coverage of the writer's whole intended message.
Given a writer's own keywords, optional cue and intended sentence, a draft set \emph{covers} the item when the writer marks some draft as the intended message or as the intended message with a small edit; the burden is the number of cue taps, drafts read and selections.
Uncovered items are attributed to loss (content the keywords could not carry), violation (a commitment the draft got wrong) or assumption (a reading the model chose without a basis), using the reason classifier of Section~\ref{sec:waterfall}.
The baselines are the single draft (intended-version match \jmpOneEightMatchBare{}\% bare and \jmpOneEightMatchChip{}\% with a chip; accepted as written \jmpOneEightAcceptYesK{} of \jmpOneEightAcceptN{}), diverse sampling, prompted diversity, and a random-dimension set that tests whether targeting the contrast, not variety, drives any gain.
The pilot design was written before any draft set was generated, and its allowed conclusions are declared in advance, including that sets help only when the cue is missing, and that they do not help beyond one cue tap; a null is reportable.
It will separate loss that no draft could recover (the keywords do not contain it) from loss that a better set could, which is the ceiling that Section~\ref{sec:waterfall} motivates.
Writers in the pilot are again classmates rating sets built from their own pairs, so its outcome is writer-intent agreement and not validation by AAC users.

\paragraph{The commitment model and its detectors.}
The whole-message result asks for an eighth dimension, event and role structure, that has an annotation scheme, a reliable detector and a human agreement estimate, none of which exist yet; the seven dimensions measured here are a floor.
Prediction-powered inference~\cite{angelopoulos2023ppi} offers a route to correct machine-judge estimates with a small human sample, which we regard as a methods addition, not a replacement for the human arbiter.

\paragraph{Boundary conditions.}
Machine-only explorations suggest that a conversation partner's preceding remark and the language of the message each move the default; these call for human-labelled checks before any claim.
In particular, whether a synthetic voice carries a contrast that the text carries (for example certainty or politeness) is a separate question that no part of this study addresses.
